\chapter{亀座標}
\label{chap:tortoise}
\chapterepigraph{``Put off thy shoes from off thy feet, for the place whereon thou standest is holy ground.''}{Exodus 3:5, King James Version}

\section{ブラックホールは開いた波動系である}

力学において一般化座標を選ぶ目的は、同じ運動を別の記号で書くことではない。
拘束や対称性を座標そのものに反映させることにある。\term{ブラックホール摂動論}で
同じ役割を果たすのが亀座標である。

遠方からブラックホールへ波束を送る。一部は\term{有効ポテンシャル}で反射され、
一部は\term{地平面}を通過する。一度\term{地平面}へ達した信号は外部領域へ戻らない。したがって、
外部領域だけを系とみなすと、エネルギーは二つの\term{漸近端}へ流出する。これは箱の中の
定常波ではなく、二端点をもつ\term{散乱問題}である。

本章では\term{幾何学単位系}
\begin{equation}
G=c=1
\label{eq:geometrized_units}
\end{equation}
を用い、計量の符号を $(-,+,+,+)$ とする。$G$ は Newton 定数、$c$ は光速度で
ある。必要になる背景時空は、まず静的かつ球対称なものに限る。

\section{静的球対称時空}

外部領域の計量を
\begin{equation}
\rmd s^2
=
-f(r)\,\rmd t^2
+\frac{\rmd r^2}{f(r)}
+r^2\rmd\Omega^2
\label{eq:sss_metric}
\end{equation}
と書く。$t$ は\term{静的時刻}、$r$ は\term{面積半径}であり、
\begin{equation}
\rmd\Omega^2
=
\rmd\vartheta^2
+\sin^2\vartheta\,\rmd\varphi^2
\end{equation}
は単位二次元球面の計量である。関数 $f(r)$ が単純零点 $r=r_h$ をもつとする。
すなわち、
\begin{align}
f(r_h)&=0,
&
f'(r_h)&\ne 0,
\label{eq:simple_horizon}
\end{align}
である。プライムは $r$ による微分を表す。$r_h$ は \term{Killing 地平面}の半径である。

計量\eqref{eq:sss_metric}では $g_{rr}=f^{-1}$ が地平面で発散する。しかし、これは
曲率が発散するという意味ではない。静的座標が、地平面を横切る光線の記述に適して
いないのである。波動方程式の境界条件を定める前に、この座標上の見かけを取り除く
必要がある。

\section{動径光線から亀座標を導く}

角度を固定した\term{動径零測地線}では、$\rmd\Omega=0$ および $\rmd s^2=0$ である。
したがって、式\eqref{eq:sss_metric}から
\begin{equation}
\frac{\rmd t}{\rmd r}
=
\mathord{\pm}\frac{1}{f(r)}
\label{eq:null_slope_r}
\end{equation}
を得る。ここで、動径座標 $r$ の代わりに $\rstar$ を
\begin{equation}
\frac{\rmd \rstar}{\rmd r}
=
\frac{1}{f(r)}
\label{eq:tortoise_definition}
\end{equation}
によって定義する。$\rstar$ を\term{亀座標}とよぶ。加法定数は任意であり、散乱振幅の
位相規約に対応する。

式\eqref{eq:null_slope_r}は
\begin{equation}
\frac{\rmd t}{\rmd \rstar}
=
\mathord{\pm}1
\end{equation}
となる。したがって、
\begin{align}
u&=t-\rstar ,
&
v&=t+\rstar
\label{eq:null_coordinates}
\end{align}
をそれぞれ\term{遅延時刻}、\term{先進時刻}と定義すれば、外向き光線では $u$ が一定、内向き
光線では $v$ が一定である。

地平面の近傍で $f(r)$ を一次まで展開すると、
\begin{equation}
f(r)
=
f'(r_h)(r-r_h)
+O\bigl((r-r_h)^2\bigr)
\end{equation}
である。式\eqref{eq:tortoise_definition}を積分して、
\begin{equation}
\rstar
=
\frac{1}{f'(r_h)}
\log|r-r_h|
+O(1)
\label{eq:tortoise_near_horizon}
\end{equation}
を得る。単純な\term{事象地平面}は、外部領域から見れば $\rstar=-\infty$ に移る。
亀座標は地平面を除去したのではなく、波の境界条件を課すべき\term{漸近端}へ移したのである。

\begin{principlebox}
亀座標は、地平面の\term{因果構造}を一次元散乱の漸近構造へ翻訳する。地平面が
$\rstar=-\infty$ に置かれることによって、「地平面へ入る波」は\term{平面波}の向きとして
定義できる。
\end{principlebox}

\section{地平面を横切る座標}

亀座標は地平面を無限遠へ送るが、地平面そのものを横切るには先進時刻
$v=t+\rstar$ を座標として用いる。$\rmd t=\rmd v-f^{-1}\rmd r$ を
式\eqref{eq:sss_metric}へ代入すると、
\begin{equation}
\rmd s^2
=
-f(r)\,\rmd v^2
+2\,\rmd v\rmd r
+r^2\rmd\Omega^2
\label{eq:ingoing-ef-metric}
\end{equation}
となる。これは $f(r_h)=0$ でも退化せず、未来地平面を横切る内向き
\term{Eddington--Finkelstein 座標}である。静的座標における $g_{rr}$ の発散が座標の見かけ
であったことも、この式から分かる。

時間因子を $\rme^{-\rmi\omega t}$ とすると、地平面へ入る波は
\begin{equation}
\rme^{-\rmi\omega t}\rme^{-\rmi\omega\rstar}
=
\rme^{-\rmi\omega v}
\label{eq:ingoing-ef-regular-mode}
\end{equation}
であり、正則な座標 $v$ だけの関数である。これに対して外向き波は
$u=t-\rstar$ の関数で、未来地平面上の静的外部から出てくる。したがって「事象
地平面で内向き」という条件は、符号の暗記ではなく、未来地平面上で正則かつ因果的な
解を選ぶ条件である。第\ref{chap:horizon-thermality}章では、同じ正則性を複素座標へ
延長して熱因子を得る。

\section{Schwarzschild 時空}

質量 $M>0$ の \term{Schwarzschild ブラックホール}では、
\begin{equation}
f(r)
=
1-\frac{2M}{r}
\label{eq:schwarzschild_f}
\end{equation}
であり、\term{事象地平面}は $r_h=2M$ にある。式\eqref{eq:tortoise_definition}を積分すると、
\begin{equation}
\rstar
=
r
+2M\log\left(\frac{r}{2M}-1\right)
\label{eq:schwarzschild_tortoise}
\end{equation}
を得る。加法定数は、対数の中を無次元化する規約に含めた。外部領域 $2M<r<\infty$
は、
\begin{equation}
-\infty<\rstar<\infty
\end{equation}
へ写る。

地平面近傍と無限遠での漸近形は、それぞれ
\begin{align}
\rstar
&=
2M\log\left(\frac{r-2M}{2M}\right)
+O(1),
&& r\longrightarrow 2M+0,
\label{eq:schwarzschild_tortoise_hor}
\\
\rstar
&=
r+2M\log\left(\frac{r}{2M}\right)
+O(r^{-1}),
&& r\longrightarrow\infty
\label{eq:schwarzschild_tortoise_inf}
\end{align}
である。無限遠でも $\rstar$ と $r$ は対数項だけずれる。この長距離の位相は、精密な
散乱振幅や漸近展開では無視できない。

$r$ 座標では地平面が有限位置にあるが、$\rstar$ 座標では外部領域全体が実直線に
なる。事象地平面と\term{空間的無限遠}は、それぞれ左端と右端の漸近領域である。

\section{スカラー波動方程式の一次元化}

背景計量上の\term{質量ゼロスカラー場}を $\Phi$ とし、外部源を $J$ とする。場の方程式を
\begin{equation}
\nabla_\mu\nabla^\mu\Phi
=
J
\label{eq:scalar_wave_covariant}
\end{equation}
とする。ギリシャ添字は四次元時空座標を走る。\term{球面調和関数} $Y_{\ell m}$ を用いて、
\begin{equation}
\Phi(t,r,\vartheta,\varphi)
=
\sum_{\ell=0}^{\infty}
\sum_{m=-\ell}^{\ell}
\frac{\Psi_{\ell m}(t,r)}{r}
Y_{\ell m}(\vartheta,\varphi)
\label{eq:scalar_decomposition}
\end{equation}
と展開する。整数 $\ell$ と $m$ は\term{角運動量量子数}であり、
\begin{equation}
\Delta_{S^2}Y_{\ell m}
=
-\ell(\ell+1)Y_{\ell m}
\end{equation}
を満たす。$\Delta_{S^2}$ は単位球面上の \term{Laplace 演算子}である。$1/r$ をあらかじめ
取り出した理由は、動径方程式から一階微分項を消すためである。

同様に源を球面調和展開して式\eqref{eq:scalar_wave_covariant}へ代入すると、各
$(\ell,m)$ に対して
\begin{equation}
\left[
\frac{\partial^2}{\partial t^2}
-\frac{\partial^2}{\partial \rstar^2}
+V_\ell(r)
\right]
\Psi_{\ell m}(t,r)
=
S_{\ell m}(t,r)
\label{eq:master_wave_time}
\end{equation}
を得る。$S_{\ell m}$ は $J$ から定まる一次元の源であり、\term{有効ポテンシャル}は
\begin{equation}
V_\ell(r)
=
f(r)
\left[
\frac{\ell(\ell+1)}{r^2}
+\frac{f'(r)}{r}
\right]
\label{eq:scalar_effective_potential}
\end{equation}
である。

\paragraph{一階微分項が消える理由.}
動径部分に現れる組合せは
\begin{equation}
\frac{1}{r^2}
\frac{\partial}{\partial r}
\left(
r^2f(r)
\frac{\partial}{\partial r}
\frac{\Psi}{r}
\right)
\end{equation}
である。$\partial_r=f^{-1}\partial_{\rstar}$ を用いて展開すると、$\Psi/r$ の微分から
生じる項と体積要素 $r^2$ の微分から生じる項が組み合わさる。その結果、
$\partial_{\rstar}\Psi$ に比例する項が消え、式\eqref{eq:master_wave_time}の
Schr\"odinger 型が得られる。場の再定義 $\Phi=\Psi/r$ は単なる慣習ではなく、
動径演算子を対称な形へ移す変換である。

Schwarzschild 時空では、式\eqref{eq:schwarzschild_f}を用いて
\begin{equation}
V_\ell(r)
=
\left(1-\frac{2M}{r}\right)
\left[
\frac{\ell(\ell+1)}{r^2}
+\frac{2M}{r^3}
\right]
\label{eq:schwarzschild_scalar_potential}
\end{equation}
となる。このポテンシャルは地平面と無限遠の両方で零に近づき、その中間に障壁を
もつ。したがって、外部摂動は一次元の障壁散乱として理解できる。

\section{時間規約と放射条件}

周波数領域へ移るため、時間依存性を
\begin{equation}
\Psi_{\ell m}(t,r)
=
\rme^{-\rmi\omega t}\psi_{\ell m}(r;\omega)
\label{eq:time_convention}
\end{equation}
と定める。$\omega\in\bbC$ は\term{複素周波数}である。ポテンシャルが消える漸近領域では、
動径関数は
\begin{equation}
\psi(\rstar;\omega)
\sim
\rme^{\mathord{\pm}\rmi\omega \rstar}
\label{eq:asymptotic_plane_waves}
\end{equation}
となる。式\eqref{eq:time_convention}と組み合わせると、
\begin{align}
\rme^{-\rmi\omega t}\rme^{-\rmi\omega \rstar}
&=
\rme^{-\rmi\omega v},
\label{eq:ingoing_wave}
\\
\rme^{-\rmi\omega t}\rme^{+\rmi\omega \rstar}
&=
\rme^{-\rmi\omega u}.
\label{eq:outgoing_wave}
\end{align}
先進時刻 $v$ の関数は内向き、遅延時刻 $u$ の関数は外向きである。したがって、
事象地平面で純粋に内向き、無限遠で純粋に外向きという\term{放射条件}は
\begin{align}
\psi(\rstar;\omega)
&\sim
\rme^{-\rmi\omega \rstar},
&& \rstar\longrightarrow-\infty,
\label{eq:bh_ingoing_bc}
\\
\psi(\rstar;\omega)
&\sim
\rme^{+\rmi\omega \rstar},
&& \rstar\longrightarrow+\infty
\label{eq:infinity_outgoing_bc}
\end{align}
となる。

\begin{warningbox}
「$\rme^{+\rmi\omega \rstar}$ が出射波」という判定は、それだけでは意味をもたない。
時間因子 $\rme^{-\rmi\omega t}$ と座標の向きを同時に指定して初めて、入射と出射が
決まる。時間規約を $\rme^{+\rmi\omega t}$ に変えれば、すべての符号が反転する。
\end{warningbox}

$\im\omega<0$ の\term{準固有振動}（quasinormal mode; \term{QNM}）周波数\footnote{%
準固有振動の定義と三つの同値な特徴づけは第\ref{chap:qnm}章で与える。}
を想定すると、式\eqref{eq:bh_ingoing_bc}と
\eqref{eq:infinity_outgoing_bc}の空間波形は実軸上で発散する。この事実は異常ではない。
時間的に減衰しながら、波面が無限遠へ進む開放系の共鳴状態を、全空間で同じ複素
周波数により表した結果である。ただし、この発散のため QNM 波形は通常の
\term{自乗可積分}な $L^2$ 固有関数にはならない。これが解析接続と、第
\ref{chap:complex-scaling}章で定義する複素スケーリングを必要とする第一の理由である。

\section{二つの地平面をもつ場合}

\term{Schwarzschild--de~Sitter 時空}では、静的領域が\term{ブラックホール地平面} $r_b$ と
\term{宇宙論的地平面} $r_c$ の間にある。ここで
\begin{equation}
r_b<r<r_c
\end{equation}
である。適切な加法定数を選ぶと、
\begin{align}
\rstar&\longrightarrow-\infty,
&&r\longrightarrow r_b+0,
\\
\rstar&\longrightarrow+\infty,
&&r\longrightarrow r_c-0
\end{align}
となる。QNM の\term{放射条件}は、ブラックホール地平面へ入る波と、宇宙論的地平面へ
出ていく波である。式の形は\eqref{eq:bh_ingoing_bc}、
\eqref{eq:infinity_outgoing_bc}と同じでも、右端の物理的意味は空間的無限遠ではなく
宇宙論的地平面になる。

一方、\term{極限ブラックホール}では $f(r)$ の零点が重なるため、
式\eqref{eq:tortoise_near_horizon}の対数則が変わる。重根 $r=r_e$ の近傍で
\begin{equation}
f(r)
=
a(r-r_e)^2
+O\bigl((r-r_e)^3\bigr)
\end{equation}
ならば、$a$ を二次項の係数として
\begin{equation}
\rstar
=
-\frac{1}{a(r-r_e)}
+O\bigl(\log|r-r_e|\bigr)
\end{equation}
となる。非極限系の公式でパラメータを極限値へ代入するだけでは、漸近構造を取り違える
ことがある。

\section{本章の結論}

本章で得た方程式は
\begin{equation}
\left(
\partial_t^2
-\partial_{\rstar}^2
+V
\right)\Psi
=
S
\label{eq:generic_master_equation}
\end{equation}
である。ブラックホールの種類と摂動のスピンは有効ポテンシャル $V$ に入り、因果構造は
$\rstar$ の範囲と漸近条件に入る。この分離によって、幾何学の問題は散乱理論の問題へ
移された。次章では、式\eqref{eq:generic_master_equation}を外力に対する応答として解く。

\begin{researchproblem}[極限地平面の散乱座標]
\term{Reissner--Nordstr\"om 時空}の
$f(r)=1-2M/r+Q^2/r^2$ について、非極限と極限の亀座標を別々に構成せよ。
$|Q|\to M$ と $r\to r_+$ の順序を交換し、近地平面の波動方程式と \term{Jost 規格化}が
どこで非一様になるかを決定せよ。既知の式を再現するだけでなく、後の数値計算で
用いる一様座標または \term{matched asymptotic expansion} を提案し、その誤差を評価せよ。
\end{researchproblem}
